\chapter{Sécurité}\label{chap:securite}

La sécurité est intégrée à la structure de l'application. Le tableau~\ref{tab:menaces} associe chaque menace courante
\cite{owaspcsrf,owasppwd} à la protection mise en place et à l'endroit où elle se trouve.

\begin{table}[htbp]
\centering
\caption{Menaces et protections}\label{tab:menaces}
\begin{tabularx}{\linewidth}{@{}lXl@{}}
\toprule
\textbf{Menace} & \textbf{Protection} & \textbf{Où}\\
\midrule
Injection SQL & Requêtes JPQL paramétrées ; saisies de recherche échappées pour \code{LIKE} & services\\
XSS & Échappement de toute sortie (\code{c:out}) ; JSON échappé ; politique CSP sans script en ligne & vues, filtre\\
CSRF & Jeton de session obligatoire sur tout POST (champ \code{\_csrf} ou en-tête \code{X-CSRF-Token}) & \code{SecurityFilter}\\
Fixation de session & Changement d'identifiant de session et nouveau jeton CSRF à la connexion & \code{ClientServlet}\\
Mots de passe volés & Hachage PBKDF2-HMAC-SHA256, sel aléatoire, 600\,000 itérations & \code{PasswordHasher}\\
Énumération de comptes & Même message pour « e-mail inconnu » et « mauvais mot de passe » ; hachage factice quand le compte n'existe pas & \code{JpaInternauteService}\\
Force brute & 5 échecs en 15 minutes par couple adresse IP + e-mail, puis blocage temporaire & \code{LoginThrottle}\\
Redirection ouverte & Le retour après connexion n'accepte qu'un chemin interne de l'application & \code{BaseServlet}\\
Accès non autorisé & Contrôle par adresse dans un filtre (\code{/panier}, \code{/admin/*}) & \code{AuthFilter}\\
Clickjacking & \code{X-Frame-Options: DENY} et \code{frame-ancestors 'none'} & \code{SecurityFilter}\\
Fuite d'informations & Pages d'erreur génériques ; pages sensibles en \code{Cache-Control: no-store} & \code{web.xml}, filtre\\
\bottomrule
\end{tabularx}
\end{table}

\section{Jeton CSRF}
Un site malveillant peut faire envoyer au navigateur d'un client connecté une requête POST vers l'application, avec son
cookie de session. Sans défense, l'action s'exécuterait. Le jeton, propre à la session, est inséré dans chaque
formulaire et comparé à celui de la session à l'arrivée du POST ; le site malveillant ne peut pas le connaître. Il est changé à la connexion.
La comparaison se fait en temps constant, pour ne rien révéler par la durée du calcul.

\section{Mots de passe}
Un mot de passe n'est jamais stocké. La base contient une empreinte au format
\code{pbkdf2\_sha256\$itérations\$sel\$hachage}. PBKDF2 \cite{rfc8018} est volontairement lent : un attaquant qui
volerait la base devrait payer ce coût pour chaque essai. Le nombre d'itérations est écrit dans l'empreinte, ce qui permet de
l'augmenter plus tard sans invalider les comptes existants.

\section{En-têtes et politique CSP}
La politique de sécurité du contenu (CSP) limite l'origine des ressources : scripts uniquement depuis l'application,
aucun script en ligne, formulaires envoyés uniquement vers l'application, pas d'intégration dans une autre page. C'est
la raison pour laquelle le JavaScript est dans un fichier séparé et non dans les pages. Les autres en-têtes sont
\code{X-Content-Type-Options: nosniff} et \code{Referrer-Policy: same-origin}. Le cookie de session est \code{HttpOnly} et
la session expire après 30 minutes d'inactivité.

\section{Ce que la sécurité ne couvre pas}
Ces limites sont assumées et reprises au chapitre~\ref{chap:perspectives} : pas de HTTPS (donc pas d'attribut \code{Secure} sur le
cookie) car l'application tourne en local ; pas de récupération de mot de passe oublié ; le blocage de la force brute est
conservé en mémoire et disparaît au redémarrage ; la politique CSP autorise encore les styles en ligne.
